% GENERATED FILE. Edit canonical lessons, then run npm run generate:books.
% canonical-source-hash: fnv1a64:283cc4cdb488731d
% canonical-lessons: KA-C52-door, KA-C52-chair, KA-C52-rangoli, KA-C52-flower, KA-C52-garland

\chapter{Welcoming a Guest}
\label{ch:ka-welcome}
\hlchaptermodality{52}

\begin{chapteropening}
\textbf{By the end of this chapter:} \emph{I can name the five things a Kannada welcome puts between the street and the guest, from the pattern on the ground to the garland.}

Complete KA-C52-garland and say all five words in order.
\end{chapteropening}

\section[bāgilu]{\kn{ಬಾಗಿಲು} (bāgilu) — a door}

\label{lesson:KA-C52-door}



\begin{quote}
Before the new run: what did \emph{vandane} mean, and what did \emph{āśīrvāda} mean?
\end{quote}

\subsection*{You'll want to know: \kn{ಬಾಗಿಲು}}
\textbf{\kn{ಬಾಗಿಲು}} (\emph{bāgilu}) — \textquotedblleft{}a door\textquotedblright{}.

\kn{ಬಾಗಿಲು} is the door and the doorway at once — the opening quite as much as the plank that shuts it.

Beside it sits Tamil \ta{வாயில்} (\emph{vāyil}), a gateway, and that word is transparently built on \ta{வாய்}, 'mouth', which is \kn{ಬಾಯಿ} here with the \emph{b} standing in for the \emph{v}. Whether \kn{ಬಾಗಿಲು} is that same word with a consonant grown in its middle, or a separate formation on \kn{ಬಾಗು}, 'to bend', is not settled. What is beyond doubt is that a Dravidian house has been known to name its opening after a face.

\begin{grammarlens}[title={What you've built}]
The first of the five things a welcome puts in a guest's way.
\end{grammarlens}

\subsection*{Your turn}
Say these aloud:

\begin{itemize}
\raggedright
  \item \emph{bāgilu}
  \item it once more, slowly
  \item \emph{bāgilu}, then \emph{bāyi}, and hear the \emph{b} and the \emph{v} trade places
\end{itemize}

\begin{tcolorbox}[breakable,colback=teal!4,colframe=teal!35!black,title={Before you move on}]
What does \kn{ಬಾಗಿಲು} mean? (\textquotedblleft{}a door\textquotedblright{}.) Say it once more.
\end{tcolorbox}
\section[kurci]{\kn{ಕುರ್ಚಿ} (kurci) — a chair}

\label{lesson:KA-C52-chair}



\begin{quote}
Before the new one: what did \emph{bāgilu} mean?
\end{quote}

\subsection*{You'll want to know: \kn{ಕುರ್ಚಿ}}
\textbf{\kn{ಕುರ್ಚಿ}} (\emph{kurci}) — \textquotedblleft{}a chair\textquotedblright{}.

Arabic \emph{kursī}, a seat or a throne, carried into India by Persian and handed on to half the languages of the subcontinent — Hindi \dv{कुर्सी}, Marathi \dv{खुर्ची}, and this.

The word was borrowed because the object was. A Kannada household ate, and to a fair degree still eats, sitting on the floor; the chair came in as furniture and brought its name along. You already know how to tell a guest to sit. Now there is something to sit on.

\begin{grammarlens}[title={What you've built}]
Two: a door and a chair.
\end{grammarlens}

\subsection*{Your turn}
Say these aloud:

\begin{itemize}
\raggedright
  \item \emph{kurci}
  \item it once more, slowly
  \item \emph{kuḷitukō}, then \emph{kurci} — the command, then the thing it needs
\end{itemize}

\begin{tcolorbox}[breakable,colback=teal!4,colframe=teal!35!black,title={Before you move on}]
What does \kn{ಕುರ್ಚಿ} mean? (\textquotedblleft{}a chair\textquotedblright{}.) Say it once more.
\end{tcolorbox}
\section[raṅgōli]{\kn{ರಂಗೋಲಿ} (raṅgōli) — the pattern drawn at a doorway}

\label{lesson:KA-C52-rangoli}



\begin{quote}
Before the new one: what did \emph{kurci} mean?
\end{quote}

\subsection*{You'll want to know: \kn{ರಂಗೋಲಿ}}
\textbf{\kn{ರಂಗೋಲಿ}} (\emph{raṅgōli}) — \textquotedblleft{}the pattern drawn at a doorway\textquotedblright{}.

White rice-flour, drawn freehand on swept ground in front of the door, at first light, before anyone has arrived.

The name is \emph{raṅga}, 'colour', and \emph{āvalī}, 'a row': a row of colour. That \emph{āvalī} is the same one behind the festival of lamps whose name came apart a few chapters back — a row of lamps there, a row of colour here. Kannada also calls the thing \kn{ರಂಗವಲ್ಲಿ} (\emph{raṅgavalli}), which is the same compound less worn down.

\begin{grammarlens}[title={What you've built}]
Three: a door, a chair, and the pattern on the ground before both.
\end{grammarlens}

\subsection*{Your turn}
Say these aloud:

\begin{itemize}
\raggedright
  \item \emph{raṅgōli}
  \item it once more, slowly
  \item \emph{raṅgōli}, then \emph{dīpa}, and name the row inside each
\end{itemize}

\begin{tcolorbox}[breakable,colback=teal!4,colframe=teal!35!black,title={Before you move on}]
What does \kn{ರಂಗೋಲಿ} mean? (\textquotedblleft{}the pattern drawn at a doorway\textquotedblright{}.) Say it once more.
\end{tcolorbox}
\section[hūvu]{\kn{ಹೂವು} (hūvu) — a flower}

\label{lesson:KA-C52-flower}



\begin{quote}
Before the new one: what did \emph{raṅgōli} mean?
\end{quote}

\subsection*{You'll want to know: \kn{ಹೂವು}}
\textbf{\kn{ಹೂವು}} (\emph{hūvu}) — \textquotedblleft{}a flower\textquotedblright{}.

Proto-Dravidian \textbf{*pū}, one syllable, and one syllable is what it stayed nearly everywhere: Tamil \ta{பூ}, Malayalam \ml{പൂവ്}, Telugu \te{పువ్వు}.

Kannada's \kn{ಹೂವು} shows the law at its very clearest. There is almost nothing else in the word — no root to hide behind, no ending to distract you — only the one sound, \emph{p} across the family and \emph{h} here.

\begin{grammarlens}[title={What you've built}]
Four. A door, a chair, the pattern, a flower.
\end{grammarlens}

\subsection*{Your turn}
Say these aloud:

\begin{itemize}
\raggedright
  \item \emph{hūvu}
  \item it once more, slowly
  \item \emph{hūvu}, then \emph{haṇṇu} and \emph{hallu}, three inherited words with one \emph{h} each
\end{itemize}

\begin{tcolorbox}[breakable,colback=teal!4,colframe=teal!35!black,title={Before you move on}]
What does \kn{ಹೂವು} mean? (\textquotedblleft{}a flower\textquotedblright{}.) Say it once more.
\end{tcolorbox}
\section[hūmāle]{\kn{ಹೂಮಾಲೆ} (hūmāle) — a garland}

\label{lesson:KA-C52-garland}



\begin{quote}
Before the new one: what did \emph{hūvu} mean?
\end{quote}

\subsection*{You'll want to know: \kn{ಹೂಮಾಲೆ}}
\textbf{\kn{ಹೂಮಾಲೆ}} (\emph{hūmāle}) — \textquotedblleft{}a garland\textquotedblright{}.

\kn{ಹೂ}, the flower from the lesson before, and \kn{ಮಾಲೆ} (\emph{māle}), from Sanskrit \emph{mālā}, a garland or any row of strung things.

So the last word of the run carries both halves of Kannada in it: a Dravidian syllable some thousands of years old, and a Sanskrit one that arrived later and settled in. Nearly every page here has turned on one of those two layers. This word declines to choose.

Five things a Kannada welcome puts between the street and the guest: the door, the chair, the pattern on the ground, the flower, the garland.

\begin{grammarlens}[title={What you've built}]
Five, and the last of them is half inherited and half borrowed.
\end{grammarlens}

\subsection*{Your turn}
Say these aloud:

\begin{itemize}
\raggedright
  \item \emph{hūmāle}
  \item it once more, slowly
  \item all five in order, then \emph{hūvu} and \emph{hūmāle} one after the other
\end{itemize}

\begin{tcolorbox}[breakable,colback=teal!4,colframe=teal!35!black,title={Before you move on}]
What does \kn{ಹೂಮಾಲೆ} mean? (\textquotedblleft{}a garland\textquotedblright{}.) Say it once more.
\end{tcolorbox}
